\begin{Corollary}\label{cor:GeneralizedVermaIrreducible}
	Assume that $\Im(\Phi)$ for $M = \ind{g}{q}(F)$ is a direct summand of $\ind{g}{q}(F)|_{\lie{g'}}$.
	Then the $\univ{g}^{G'}$-module $\Hom_{\lie{g'}}\bigl(\ind{g'}{q'}(F'), \ind{g}{q}(F)\bigr)$ is irreducible or zero.
\end{Corollary}

\begin{Remark}
	For example, if $\ind{g}{q}(F)|_{\lie{g'}}$ is completely reducible, the assumption of the corollary holds.
\end{Remark}

\begin{proof}[Proof of Corollary~\ref{cor:GeneralizedVermaIrreducible}]
	Let $[\lambda']$ denote the infinitesimal character of $\ind{g'}{q'}(F')$.
	By Proposition~\ref{prop:AdmissibleDual}, the primary component $P_{[\lambda']}(M|_{\lie{g'}})$ has finite length.
	Since $M|_{\lie{g'}}$ has a standard filtration, so does $P_{[\lambda']}(M|_{\lie{g'}})$.
	By~\cite[Proposition~3.7 and Theorem~9.8\,(3)]{Hu08_category_o}, any direct summand of $P_{[\lambda']}(M|_{\lie{g'}})$ has a standard filtration.
	In particular, $\Coker(\Phi)$ has a standard filtration by assumption.
	Therefore, the assertion follows from Theorems~\ref{thm:Exactness} and~\ref{thm:InjectivePhi} by letting $M'$ be an~irreducible quotient of \smash{$\ind{g'}{q'}\bigl(\fdual{F'}\bigr)$}.
	See Proposition~\ref{prop:CompletelyReducibleVerma} and Remark~\ref{rmk:InjectivePhiStandard} for the vanishing of~\smash{$\Dzuck{\Delta(G')}{\Delta(L')}{S - 1}(M'\otimes \Coker(\Phi))^{\Delta(G')}$}.
\end{proof}

\subsection{Quasi-abelian parabolic subalgebra}

In this subsection, we give a sufficient condition for completely reducibility of $\ind{g}{q}(F)|_{\lie{g'}}$
to apply Corollary~\ref{cor:GeneralizedVermaIrreducible}.
Retain the notation in the previous subsection.

Set $\lie{u}''\coloneq \lie{u} \cap (\lie{g}')^\perp$ and $\lie{\overline{u}}''\coloneq \lie{\overline{u}} \cap (\lie{g}')^\perp$.
Then $\lie{u}''$ is $\lie{u}'$-stable and $\lie{u} = \lie{u}' \oplus \lie{u}''$ holds by the definition of weakly $\lie{g'}$-compatible parabolic subalgebras (Definition~\ref{def:WeaklyCompatible}).
The notion of quasi-abelian parabolic subalgebras is defined in~\cite[p.~109]{EnPaWaWo85} for symmetric $(\lie{g}, \lie{g'})$.
We consider a~straightforward generalization of the definition.

\begin{Definition}\label{def:QuasiAbelian}
	$\lie{q}$ is said to be \emph{quasi-abelian} with respect to $\lie{g'}$ if
	$(\alpha, \beta)\geq 0$ holds for any $\alpha \in \Delta(\lie{u'}, \lie{t'})$ and $\beta \in \Delta(\lie{u''}, \lie{t'})$.
\end{Definition}

It is clear that if $\lie{p}$ is a weakly $\lie{g'}$-compatible parabolic subalgebra containing $\lie{q}$ and $\lie{q}$ is quasi-abelian, then $\lie{p}$ is also quasi-abelian.
In fact, the nilpotent radical of $\lie{p}$ is smaller than that of~$\lie{q}$.

If the nilpotent radical of $\lie{q}$ is abelian, $\lie{q}$ is quasi-abelian.
More generally, we have the following criterion.

\begin{Proposition}\label{prop:QuasiAbelianCondition}
	If $[\lie{u}', \lie{u}'']=0$ holds, then $\lie{q}$ is quasi-abelian.
	In particular, a parabolic subalgebra with abelian nilpotent radical is quasi-abelian.
\end{Proposition}
	
\begin{proof}
	Let $\alpha \in \Delta(\lie{u'}, \lie{t'})$ and $\beta \in \Delta(\lie{u''}, \lie{t'})$.
	Assume $(\alpha, \beta) < 0$.
	This implies $\alpha + \beta \in \Delta(\lie{u''}, \lie{t'})$
	and \smash{$\bigl[\lie{u}'_{\alpha}, \lie{u}''_{\beta}\bigr]\neq 0$}.
	By assumption, \smash{$\bigl[\lie{u}'_{\alpha}, \lie{u}''_{\beta}\bigr]\subset [\lie{u}', \lie{u}'']=0$} holds and this is contradiction.
	Therefore, $\lie{q}$ is quasi-abelian.
\end{proof}

\begin{Example}
	If $\lie{g'} \simeq \lie{sl}(2,\CC)$, then $\lie{q}$ is quasi-abelian as follows.
	$\lie{q}$ contains $\lie{q}(H)$ defined by a semisimple element $H\in \lie{g'}$.
	Since any eigenvalue in $\lie{u}(H)$ of $\ad(H)$ is positive, $\lie{q}(H)$ is quasi-abelian.
	As noted immediately after Definition~\ref{def:QuasiAbelian}, $\lie{q}$ is also quasi-abelian.

	Take an $\lie{sl}_2$-triple $\set{H, X, Y}\subset \lie{g'}$ with $[X, Y] = H$.
	If $\ad(H)$ has an eigenvalue greater than or equal to $3$ in $\lie{g}$, then we have $[\lie{u'}, \lie{u}''] \neq 0$
	by the representation theory of $\lie{sl}(2,\CC)$.
	Take a~finite-dimensional irreducible $\lie{sl}(2,\CC)$-module $F$ of dimension greater than or equal to $3$.
	Then the embedding $\lie{g'} = \lie{sl}(2,\CC) \hookrightarrow \lie{sl}(F) = \lie{g}$ gives an example
	of quasi-abelian parabolic subalgebra with $[\lie{u'}, \lie{u''}] \neq 0$.
	This means that the converse of Proposition~\ref{prop:QuasiAbelianCondition} does not hold in general.
\end{Example}


A criterion for the completely reducibility of $\ind{g}{q}(F)|_{\lie{g'}}$ is given in~\cite[Lemma~3.1]{EnPaWaWo85} if $(\lie{g}, \lie{g'})$ is a symmetric pair.
We extend the criterion to our setting.

\begin{Lemma}\label{lem:QuasiAbelianCompletelyReducible}
	Let $F$ be a finite-dimensional irreducible $\lie{l}$-module.
	Assume that for any irreducible submodule of $F|_{\lie{l}'}$, its infinitesimal character $[\lambda'] \in (\lie{t'})^*/W_{L'}$ satisfies
	\begin{align*}
		\frac{2(\lambda'+\rho(\lie{u}'), \alpha)}{(\alpha, \alpha)} \not \in \set{1, 2, \ldots}, \qquad \forall \alpha \in \Delta(\lie{u'}, \lie{t'}).
	\end{align*}
	Suppose that $\lie{q}$ is quasi-abelian.
	Then \smash{$\ind{g}{q}(F)|_{\lie{g}'}$} is completely reducible and each irreducible component is isomorphic to an irreducible generalized Verma module.
\end{Lemma}

\begin{proof}
	By Proposition~\ref{prop:FiltrationRestriction}, \smash{$\ind{g}{q}(F)|_{\lie{g}'}$} has a standard filtration with
	\begin{align*}
		\gr\bigl(\ind{g}{q}(F)|_{\lie{g}'}\bigr) \simeq \ind{g'}{q'}\bigl(F\otimes S\bigl(\lie{\overline{u}}''\bigr)\bigr),
	\end{align*}
	where the $\lie{u}'$-action on $S(\lie{\overline{u}}'')$ is trivial.

	Let $V$ be an irreducible $\lie{l'}$-submodule of $F\otimes S(\lie{\overline{u}}'')$.
	By the Weyl character formula, the infinitesimal character of $V$ is of the form $[\lambda' - R]$,
	where $R$ is a sum of elements in $\Delta(\lie{u''}, \lie{t'})$ and~$[\lambda']$ is the infinitesimal character of some irreducible submodule of $F|_{\lie{l'}}$.
	Since $\lie{q}$ is quasi-abelian, we have
	\begin{align*}
		\frac{2(R, \alpha)}{(\alpha, \alpha)} \in \set{0, 1, 2, \ldots}, \qquad \forall \alpha \in \Delta(\lie{u'}, \lie{t'}).
	\end{align*}
	This and the assumption imply
	\begin{align*}
		\frac{2(\lambda'-R+\rho(\lie{u}'), \alpha)}{(\alpha, \alpha)} \not \in \set{1, 2, \ldots}\qquad \text{for any }\alpha \in \Delta(\lie{u'}, \lie{t'}).
	\end{align*}
	Hence \smash{$\ind{g'}{q'}(V)$} is irreducible by Fact~\ref{fact:VermaIrreducible}.
	Therefore, Proposition~\ref{prop:CompletelyReducibleVerma} shows the assertion.
\end{proof}

The condition in Lemma~\ref{lem:QuasiAbelianCompletelyReducible} is given in terms of $\lie{t'}$-weights.
We shall give a criterion for the completely reducibility in terms of $\lie{t}$-weights.
To do so, we prepare several convexity results.
Note that J.A.~Vargas gave a similar estimate in~\cite[\S(1.5)]{Va16} in the context of the branching problem of discrete series representations.
Our result (Theorem~\ref{thm:QuasiAbelianGoodRange}) contains his estimate.
See also Theorems~\ref{thm:DerivedFunctorModuleQuasiAbelian} and~\ref{thm:DiscreteSeries}.

Suppose $\lie{t'}\subset \lie{t}$.
Let $\lie{t}''$ be the orthogonal complement of $\lie{t}'$ in $\lie{t}$.
Consider $(\lie{t'})^*$ as a subspace of $\lie{t}^*$
using the direct sum decomposition $\lie{t}=\lie{t'} \oplus \lie{t''}$.
For $\alpha \in \Delta(\lie{u'}, \lie{t'})$, define
\begin{align*}
	\Delta(\alpha)\coloneq \set{\beta \in \Delta(\lie{u}, \lie{t})\mid \beta|_{\lie{t}'}=\alpha}.
\end{align*}
For a subset $S$ of a real vector space, we denote by $\Co(S)$ the convex hull of $S$.

\begin{Lemma}\label{lem:ConvexityRestriction}
	Let $\alpha \in \Delta(\lie{u'}, \lie{t'})$.
	Then $\alpha \in \Co(\Delta(\alpha))$ holds.
\end{Lemma}

\begin{proof}
	Fix a Cartan involution $\theta$ of $\lie{g}$ such that
	$\lie{g}', \lie{t}$ and $\lie{t}'$ are $\theta$-stable.
	In fact, since $\lie{g}'$ is reductive in $\lie{g}$, such an involution exists.

	Take a root vector $X \in \lie{u}'_{\alpha}$.
	By $X \in \lie{u}$, we can write \smash{$X=\sum_{\beta \in \Delta(\alpha)}X_{\beta}$} with $X_{\beta} \in \lie{u}_{\beta}$.
	Consider the inner product $\langle \cdot, \cdot \rangle \coloneq -(\cdot, \theta(\cdot))$ on $\lie{g}$.
	Note that $[\lie{t''}, \lie{g'}] \subset (\lie{g}')^\perp$.
	Then the root spaces are mutually orthogonal, and hence we have
	\begin{align*}
		&\|X\|^2= \sum_{\beta \in \Delta(\alpha)} \|X_\beta\|^2, \\
		&\sum_{\beta \in \Delta(\alpha)} \langle \beta(Z')X_\beta, X_\beta\rangle= \langle \alpha(Z')X, X\rangle, \\
		&\sum_{\beta \in \Delta(\alpha)} \langle \beta(Z'')X_\beta, X_\beta\rangle= \langle [Z'', X], X\rangle = 0 = \langle \alpha(Z'')X, X\rangle
	\end{align*}
	for any $Z' \in \lie{t'}$ and $Z'' \in \lie{t''}$.
	This shows
	\begin{align*}
		\alpha = \sum_{\beta \in \Delta(\alpha)} \frac{\|X_\beta\|^2}{\|X\|^2} \beta, \qquad \sum_{\beta \in \Delta(\alpha)} \frac{\|X_\beta\|^2}{\|X\|^2} = 1.
	\end{align*}
	Therefore, we obtain $\alpha \in \Co(\Delta(\alpha))$.
\end{proof}


\begin{Lemma}\label{lem:SumConvex}
	Fix a Borel subalgebra $\lie{b}_L = \lie{t}\oplus \lie{u}_L$ of $\lie{l}$.
	Let $\lambda_1$ and $\lambda_2$ be dominant integral weights of $\lie{l}$.
	Then we have $\Co(W_{L}\lambda_1)+\Co(W_{L}\lambda_2) = \Co(W_{L}(\lambda_1+\lambda_2))$.
\end{Lemma}
	
\begin{proof}
	$\Co(W_{L}\lambda_1)+\Co(W_{L}\lambda_2) \supset \Co(W_{L}(\lambda_1+\lambda_2))$ is obvious.
	We shall show the converse inclusion.
	Let $s \in W_L$ and it is enough to prove $\lambda_1 +s(\lambda_2) \in \Co(W_{L}(\lambda_1+\lambda_2))$.

	For a dominant integral weight $\lambda$ of $\lie{l}$, we denote by $F(\lambda)$ the irreducible $\lie{l}$-module with the highest weight $\lambda$.
	Then there exists a unique non-zero homomorphism $p\colon F(\lambda_1) \otimes F(\lambda_2) \rightarrow F(\lambda_1 + \lambda_2)$ up to scalar.
	For each $i = 1, 2$, fix a highest weight vector $v_i$ of $F(\lambda_i)$.
	Then we have $p(v_1\otimes v_2) \neq 0$.
	
	Take a weight vector $v'_2 \in F(\lambda_2)$ with the extreme weight $s(\lambda_2)$.
	Then there exists $X \in \univ{u_{\mathit{L}}}$ such that $Xv'_2=v_2$.
	Since $v_1$ is $\lie{u}_L$-invariant, we have $X(v_1\otimes v'_2)=v_1 \otimes v_2$.
	This implies
	\begin{align*}
		Xp\bigl(v_1\otimes v'_2\bigr) = p\bigl(X\bigl(v_1\otimes v'_2\bigr)\bigr) = p(v_1\otimes v_2)\neq 0,
	\end{align*}
	and hence $p\bigl(v_1\otimes v'_2\bigr) \neq 0$.
	
	Since any weight of $F(\lambda_1+\lambda_2)$ belongs to $\Co(W_{L}(\lambda_1+\lambda_2))$,
	this shows $\lambda_1+s(\lambda_2) \in \Co(W_{L}(\lambda_1+\lambda_2))$.
	We have proved the lemma.
\end{proof}

Fix a basis $\set{v_1, \ldots, v_{r'}}$ of $\lie{t'}$ and extend it to a basis $\set{v_1, \ldots, v_r}$ of $\lie{t}$.
The bases determine lexicographical orders on $(\lie{t'})^*$ and $\lie{t}^*$.
Let $\Delta^+(\lie{l'}, \lie{t'})$ and $\Delta^+(\lie{l}, \lie{t})$ denote the sets of positive roots given by the orders.
Write $\rhog{\lie{l'}}$ (resp.\ $\rhog{\lie{l}}$) for half the sum of roots in $\Delta^+(\lie{l'}, \lie{t'})$ (resp.\ $\Delta^+(\lie{l}, \lie{t})$).
Then $\rhog{\lie{l}}|_{\lie{t'}}$ is a dominant integral weight of $\lie{l}'$.

\begin{Theorem}\label{thm:QuasiAbelianGoodRange}
	Let $F$ be a finite-dimensional irreducible $\lie{l}$-module with infinitesimal character $[\lambda] \in \lie{t}^*/W_L$
	in the good range, namely,
	\begin{align*}
		\Re(\lambda+\rho(\lie{u}), \alpha) < 0, \qquad \forall \alpha \in \Delta(\lie{u}, \lie{t}).
	\end{align*}
	Suppose that $\lie{q}$ is quasi-abelian.
	Then \smash{$\ind{g}{q}(F)|_{\lie{g}'}$} is completely reducible,
	and each irreducible direct summand is of the form \smash{$\ind{g'}{q'}(F')$} such that $F'$ is a finite-dimensional irreducible $\lie{l'}$-module in the good range.
\end{Theorem}

\begin{proof}
	Let $F'$ be an irreducible submodule of $F|_{\lie{l}'}$ and $[\lambda']$ the infinitesimal character of $F'$.
	We shall prove
	$\Re(\lambda'+\rhog{\lie{u}'}, \alpha) < 0$, $\forall \alpha \in \Delta(\lie{u'}, \lie{t'})$.
	If we prove this, the assertion follows by the same way as Lemma~\ref{lem:QuasiAbelianCompletelyReducible}.

	Replacing the representatives $\lambda$ and $\lambda'$, we may assume that $\lambda$ and $\lambda'$ are dominant with respect to $\Delta^+(\lie{l}, \lie{t})$ and $\Delta^+(\lie{l'}, \lie{t'})$, respectively.
	By Lemma~\ref{lem:SumConvex}, we have
	\begin{align*}
		\lambda' &\in \Co(W_L(\lambda - \rhog{\lie{l}}))|_{\lie{t'}} + \rhog{\lie{l'}}
		 \subset \Co(W_{L}(\lambda-\rhog{\lie{l}}))|_{\lie{t'}}+\Co(W_{L}\rhog{\lie{l}})|_{\lie{t'}} - \rhog{\lie{l}}|_{\lie{t'}} + \rhog{\lie{l'}} \\
		&= \Co(W_{L}(\lambda))|_{\lie{t'}} - \rhog{\lie{l}}|_{\lie{t'}} + \rhog{\lie{l'}}.
	\end{align*}
	From this, we write
	\begin{align*}
		\lambda' = \sum_{s \in W_L} a_s s(\lambda)|_{\lie{t'}} - \rhog{\lie{l}}|_{\lie{t'}} + \rhog{\lie{l'}}
	\end{align*}
	with $\sum_{s\in W_L}a_s = 1$ and $a_s \geq 0$.

	Let $\alpha \in \Delta(\lie{u'}, \lie{t'})$.
	By Lemma~\ref{lem:ConvexityRestriction}, $\alpha$ is written as
	\begin{align*}
		\alpha = \sum_{\beta \in \Delta(\alpha)} c_{\beta} \beta
	\end{align*}
	with $\sum_{\beta \in \Delta(\alpha)}c_{\beta} = 1$ and $c_{\beta} \geq 0$.
	Then we have
	\begin{align}
		&(\lambda'+\rhog{\lie{u}'}, \alpha) + (\rhog{\lie{u}}|_{\lie{t'}} + \rhog{\lie{l}}|_{\lie{t'}} - \rhog{\lie{l'}} - \rhog{\lie{u'}}, \alpha) \nonumber \\
		&\qquad{}=\biggl(\sum_{s \in W_L} a_s s(\lambda) + \rhog{\lie{u}}, \alpha\biggr) \nonumber \\
		&\qquad{}=\sum_{s \in W_L, \beta \in \Delta(\alpha)}a_s c_{\beta}\bigl(\lambda + \rhog{\lie{u}}, s^{-1}(\beta)\bigr). \label{eqn:PositivityTo}
	\end{align}
	By $\Delta(\alpha)\subset \Delta(\lie{u}, \lie{t})$ and the assumption, the real part of \eqref{eqn:PositivityTo} is negative.
	By Proposition~\ref{prop:PositivityRho}, we have
	\begin{align*}
		(\rhog{\lie{u}}|_{\lie{t'}} + \rhog{\lie{l}}|_{\lie{t'}} - \rhog{\lie{l'}} - \rhog{\lie{u'}}, \alpha) \geq 0.
	\end{align*}
	Therefore, we obtain $(\lambda'+\rhog{\lie{u}'}, \alpha) < 0$.
\end{proof}

Combining Corollary~\ref{cor:GeneralizedVermaIrreducible} with Theorem~\ref{thm:QuasiAbelianGoodRange}, we obtain the following.

\begin{Corollary}\label{cor:IrreducibleVermaQuasiAbelian}
	Let $F$ be a finite-dimensional irreducible $\lie{l}$-module with infinitesimal character $[\lambda] \in \lie{t}^*/W_L$
	in the good range, and $F'$ a finite-dimensional irreducible $\lie{l}'$-module.
	Suppose that~$\lie{q}$ is quasi-abelian.
	Then the $\univ{g}^{G'}$-module \smash{$\Hom_{\lie{g'}}\bigl(\ind{g'}{q'}(F'), \ind{g}{q}(F)\bigr)$} is irreducible or zero.
\end{Corollary}

The estimate in the proof of Theorem~\ref{thm:QuasiAbelianGoodRange} is independent of the assumption that $\lie{q}$ is quasi-abelian.
